\newpage
\section{量子比特映射与路由联合优化}

\subsection{方法概述}

量子线路编译中的比特映射（初始布局选择）和路由（SWAP 插入）传统上被视为两个独立阶段：先用启发式算法（如 SABRE 的 \texttt{swap\_trials}）选择初始布局，再执行路由。这种分离忽略了布局质量与路由效率之间的耦合关系——一个好的初始布局可以显著减少路由阶段所需的 SWAP 数量，从而降低噪声累积。

本工作提出将映射和路由统一到同一个强化学习 MDP 中，通过映射阶段（mapping phase）实现布局-路由联合优化。智能体在映射阶段通过虚拟 SWAP 学习初始布局，在路由阶段通过物理 SWAP 完成门调度，两个阶段共享同一个策略网络和 GNN 编码器。

\subsection{统一 MDP 设计}

\subsubsection{统一动作空间}

动作空间为 $\text{Discrete}(\text{num\_edges} + 1)$，其中：
\begin{itemize}
  \item 动作 $0 \sim \text{num\_edges} - 1$：对应物理拓扑上的耦合边 $(p, q)$，执行 SWAP 操作
  \item 动作 $\text{num\_edges}$：commit 动作，结束映射阶段并切换到路由阶段
\end{itemize}

观测向量为所有耦合边的局部特征拼接，末尾拼接映射向量（每比特的物理位置编码）、进度标量（已执行门比例）和 phase 标识符（映射阶段为 1.0，路由阶段为 0.0）。

\subsubsection{虚拟 SWAP 机制}

在映射阶段，智能体选择耦合边 $(p, q)$ 执行虚拟 SWAP（virtual SWAP）。虚拟 SWAP 仅重排逻辑比特到物理比特的映射关系：

\[
\text{mapping}[p] \leftrightarrow \text{mapping}[q]
\]

关键特性：
\begin{itemize}
  \item 不写入物理线路（\texttt{\_phys\_circuit} 不变）
  \item 不递增 SWAP 计数器（\texttt{\_swap\_counter} 不变）
  \item 不执行任何门操作
  \item 仅改变映射向量，为后续路由提供更好的初始布局
\end{itemize}

虚拟 SWAP 的距离缩短奖励与路由阶段相同：
\[
r_{\text{dist}} = -\eta \cdot \frac{d_{\text{after}} - d_{\text{before}}}{\max(d_{\text{before}}, 1)}
\]
其中 $d_{\text{before}}$ 和 $d_{\text{after}}$ 分别为虚拟 SWAP 前后 front layer 门到目标边的总距离。这使得智能体在映射阶段就能获得与路由阶段一致的距离信号，学习到什么样的布局能减少后续路由的 SWAP 需求。

\subsubsection{阶段切换与自动执行}

当智能体选择 commit 动作或映射预算耗尽时，映射阶段结束：
\begin{enumerate}
  \item 系统调用 \texttt{\_auto\_execute\_batch()}，自动执行所有前置条件已满足的门（即两个操作数物理相邻的门）
  \item 观测向量中的 phase 标识符从 1.0 变为 0.0
  \item 动作空间从 $\text{Discrete}(\text{num\_edges} + 1)$ 变为 $\text{Discrete}(\text{num\_edges})$（commit 动作不再合法）
  \item 进入路由阶段，智能体开始通过物理 SWAP 完成剩余门操作
\end{enumerate}

\subsubsection{映射预算与布局质量惩罚}

最大虚拟 SWAP 数为 $\text{mapping\_budget} = n - 1$（$n$ 为逻辑比特数），防止无限映射阶段。智能体也可以在预算内提前 commit，以平衡布局探索和路由效率。

在 commit 时刻，可选地施加布局质量惩罚：
\[
r_{\text{layout}} = -\lambda_{\text{layout}} \cdot \frac{1}{|F|} \sum_{g \in F} d(\pi(g.q_1), \pi(g.q_2))
\]
其中 $F$ 为当前 front layer，$\pi$ 为学到的映射，$d$ 为物理距离。该惩罚鼓励智能体在 commit 时已经将关键门的两个操作数放置在相邻或近距离的物理比特上。

\subsection{与传统方法的对比}

\begin{table}[htbp]
\centering
\caption{映射方法对比}
\label{tab:mapping-comparison}
\small
\begin{tabular}{l|cc}
\hline
特性 & SABRE 映射 & 本工作（联合优化） \\
\hline
映射方法 & 启发式搜索（swap\_trials） & RL 策略学习 \\
布局选择 & 多次采样取最优 & 单次前馈推理 \\
布局-路由耦合 & 独立（布局固定后路由不变） & 联合（映射阶段可探索布局-路由交互） \\
噪声感知 & 仅通过保真度评估 & Phase 2 噪声感知微调 \\
泛化性 & 需要逐拓扑调参 & 多拓扑混合训练 \\
\hline
\end{tabular}
\end{table}

传统 SABRE 方法通过 \texttt{swap\_trials} 参数控制布局采样次数：对每个候选布局执行完整路由，选择 SWAP 数最少的布局。这种方法计算开销大（$O(\text{trials} \times \text{routing\_cost})$），且布局质量仅通过 SWAP 数间接评估。

本工作的联合优化方法将布局选择和路由统一到同一个 RL 训练流程中：映射阶段的虚拟 SWAP 奖励直接反映布局对后续路由的影响，策略网络通过端到端训练学习布局-路由的联合优化策略。训练完成后，映射和路由均通过单次前馈推理完成，计算效率显著提升。

\subsection{训练策略}

映射阶段的训练集成在两阶段课程训练中：

\paragraph{Phase 1（纯路由 + 映射）。} 奖励函数包含虚拟 SWAP 的距离缩短奖励和 commit 时的布局质量惩罚。策略网络学习在映射阶段选择能减少后续路由 SWAP 需求的初始布局。

\paragraph{Phase 2（噪声感知微调）。} 在 Phase 1 检查点基础上，增加终端保真度奖励。策略网络进一步学习在噪声环境下选择能提高最终保真度的布局-路由策略。

训练采用多拓扑混合（cross\_5q、ring\_5q、ibmq\_5\_line），增强策略对不同硬件拓扑的泛化能力。映射预算为 $n - 1$，布局混合策略（\texttt{--layout-mix}）在训练中随机使用恒等映射、随机映射和 SABRE 映射作为初始状态，迫使策略学习与初始布局无关的路由能力。
